% Part: counterfactuals
% Chapter: introduction
% Section: material-conditional

\documentclass[../../../include/open-logic-section]{subfiles}

\begin{document}

\olfileid[hi]{cnt}{int}{mat}

\olsection{भौतिक सशर्त}

अंग्रेज़ी में सशर्त का सबसे सरल रूप ``यदि \dots तो \dots'' आकार का वाक्य
है, जहाँ \dots{} के स्थानों पर स्वयं वाक्य होते हैं; जैसे, ``यदि यह काम
खानसामे ने किया, तो माली निर्दोष है।'' आरंभिक तर्कशास्त्र पाठ्यक्रमों में हम
$\lif$ संयोजक से सशर्तों को संकेतित करना सीखते हैं: \dots से सूचित भागों को,
मसलन !!{formula}ों $!A$ और~$!B$ से, संकेतित करें; तब पूरे सशर्त को
$!A \lif !B$ से संकेतित किया जाता है।

संयोजक $\lif$ \emph{सत्य-फलनीय} है; अर्थात $!A \lif !B$ का
सत्य-मूल्य---$\True$ या $\False$---$!A$ और~$!B$ के सत्य-मूल्यों से निर्धारित
होता है: $!A \lif !B$ तभी सत्य है जब~$!A$ असत्य या~$!B$ सत्य हो, और अन्यथा
असत्य। सत्य-मूल्य नियतन~$\pAssign v$ के सापेक्ष हम
$\pSat{v}{!A \lif !B}$ तभी परिभाषित करते हैं जब $\pSat/{v}{!A}$ या
$\pSat{v}{!B}$। इस अर्थविज्ञान वाले संयोजक~$\lif$ को \emph{भौतिक सशर्त}
कहते हैं।

इस परिभाषा से अनेक प्रारंभिक तार्किक तथ्य मिलते हैं। सबसे पहले, भौतिक सशर्त
के लिए निगमन प्रमेय सत्य है:
\begin{align}
  \text{यदि } \Gamma, !A \Entails !B \text{ तो } \Gamma \Entails !A \lif !B
\end{align}
यह सत्य-फलनीय है: $!A \lif !B$ और $\lnot !A \lor !B$ समतुल्य हैं:
\begin{align}
  !A \lif !B & \Entails \lnot !A \lor !B\\
  \lnot !A \lor !B & \Entails !A \lif !B
  \intertext{भौतिक सशर्त अपने परिणामांश और अपने पूर्ववर्ती के निषेध, दोनों
    से निष्पन्न होता है:}
  !B & \Entails !A \lif !B\\
  \lnot !A & \Entails !A \lif !B
  \intertext{असत्य भौतिक सशर्त अपने पूर्ववर्ती और परिणामांश के निषेध के
    संयोजन के समतुल्य है: यदि $!A \lif !B$ असत्य है, तो
    $!A \land \lnot !B$ सत्य है, और विलोमतः:}
  \lnot(!A \lif !B) & \Entails !A \land \lnot !B\\
  !A \land \lnot !B & \Entails \lnot(!A \lif !B)
  \intertext{भौतिक सशर्त मोडस पोनेन्स को मान्य करता है:}
  !A, !A \lif !B & \Entails !B
  \intertext{भौतिक सशर्त समुच्चित होता है:}
  !A \lif !B,  !A \lif !C & \Entails !A \lif (!B \land !C)
  \intertext{हम पूर्ववर्ती को सदैव प्रबल कर सकते हैं; अर्थात सशर्त
    \emph{एकदिष्ट} है:}
  !A \lif !B & \Entails (!A \land !C) \lif !B
  \intertext{भौतिक सशर्त संक्रमणशील है; अर्थात शृंखला-नियम मान्य है:}
  !A \lif !B, !B \lif !C & \Entails !A \lif !C
  \intertext{भौतिक सशर्त अपने प्रतिधनात्मक के समतुल्य है:}
  !A \lif !B & \Entails \lnot !B \lif \lnot !A\\
  \lnot !B \lif \lnot !A & \Entails !A \lif !B
\end{align}

ये सभी गणितीय तर्क-विचार में उपयोगी और निर्विवाद अनुमान हैं। किंतु दार्शनिक
और भाषावैज्ञानिक साहित्य में गैर-गणितीय संदर्भों के समकक्ष अनुमानों के कथित
प्रतिदृष्टांत भरे पड़े हैं। वे संकेत करते हैं कि अंग्रेज़ी के ``यदि \dots तो
\dots'' कथनों को संकेतित करने के लिए भौतिक सशर्त~$\lif$ उपयुक्त संयोजक नहीं
है---या कम-से-कम सदैव नहीं है।

\end{document}
